\chapter[Methodology and Experiments]{Methodology and Experimental Discipline}
\label{chap:methodology}

\section{Project methodology}

The project methodology was deliberately incremental. The implementation was already large
when the project started, and changing protocol logic before understanding the real execution
path would have made the gas measurements unreliable. The work was therefore organized
around four principles:

\begin{enumerate}
  \item first reproduce the inherited implementation end-to-end;
  \item then map implementation actions to protocol steps;
  \item only then implement variants and optimizations;
  \item finally measure each variant with explicit accounting scopes.
\end{enumerate}

The first phase was diagnostic. The goal was to make the full local stack run: Hardhat node,
contract deployment scripts, Rust and WebAssembly artifacts, SQLite database, Next.js
application, Alto bundler, and optional Electron shell. This phase was not only installation
work. It established which code paths were actually used by the application and which
contracts were only present as historical or testing artifacts.

The second phase was protocol mapping. The paper describes the protocol as a sequence of
steps, while the implementation exposes functions, API routes, UI actions, and benchmark
helpers. A gas number is meaningful only if the measured code can be mapped back to a
protocol action. This mapping was especially important for the dispute phase, where a single
informal label such as ``dispute cost'' can refer to deployment, trigger, challenge rounds,
Step 8 verification, finalization, or all of them together.

The third phase was implementation. The first implementation strategy was conservative:
support variants by extending the existing contracts. This was useful for functional
validation, but it produced monolithic contracts with many branches. After the first gas
measurements and feedback from the supervisor, the implementation strategy was revised:
variant selection should happen off-chain or at deployment time, and the smart contract used
for one exchange should not carry unrelated modes. This led to the specialized final
implementation described in later chapters.

The fourth phase was measurement and consistency checking. Several early measurements were
individually correct but not directly comparable because they mixed scopes. The final
methodology therefore treats measurement scope as part of the experimental result. Every
reported number must specify which implementation was measured, which protocol steps were
included, and whether the cost is one-time infrastructure cost or marginal per-exchange
cost.

\section{Reproducibility strategy}

The reproducibility strategy has two layers. The first layer is executable reproducibility:
the repository contains scripts and tests that can recreate the local stack and rerun the
measurements. The second layer is accounting reproducibility: the report explains exactly how
to recombine the measured components.

The local stack is started with the root scripts documented in the repository. In the normal
development workflow, a Hardhat node is started first, \path{deploy-all.sh} deploys the
contract stack and writes deployment information, \path{run-alto.sh} starts the Alto bundler,
and \texttt{npm\ run\ dev} starts the Next.js application. This is the same environment used for
interactive demonstrations.

For gas measurements, the main source of reproducibility is the Hardhat test suite under
\path{src/hardhat/test/performance}. These tests deploy their own local contracts, execute
transactions, and read gas from actual transaction receipts. The tests avoid relying on
manual UI actions for the final numbers. The UI is useful to demonstrate the proof of
concept, but scripted tests are more appropriate for repeatable accounting.

The benchmarks also control cryptographic inputs whenever possible. Small and medium files
are generated deterministically inside the test code. The AES key is fixed in gas tests so
that the same circuit and commitments can be reproduced. Large-file tests create temporary
input files and validate the resulting number of ciphertext blocks and gates before running
the dispute scenario. Temporary artifacts are removed after the test when the benchmark
finishes.

This strategy does not mean that all timings are machine-independent. Gas measurements are
deterministic for a fixed compiler, contract version, and local EVM configuration. Off-chain
timings depend on the machine, file system, and runtime environment. The report therefore
separates gas results from wall-clock timings and never uses local CPU time as a substitute
for on-chain gas.

\section{Mapping code functions to protocol steps}

The implementation uses names that are convenient for Solidity and frontend code, but they
do not always match the step names used in the paper. The mapping in
Table~\ref{tab:methodology-step-mapping} is the one used throughout the measurements.

\begin{table}[htbp]
  \centering
  \begin{tabularx}{\linewidth}{>{\raggedright\arraybackslash}p{0.30\linewidth}>{\raggedright\arraybackslash}p{0.17\linewidth}>{\raggedright\arraybackslash}X}
    \toprule
    Implementation action & Protocol step & Measurement consequence \\
    \midrule
    \texttt{triggerDispute} &
    Step 6 &
    Starts the dispute path. In the measured contracts this action may also deploy or create
    the dispute account, so it must not be interpreted as a small state transition only. \\
    \texttt{respondChallenge} &
    Step 7b &
    Buyer-side answer to the current challenge. Together with \texttt{giveOpinion}, it forms
    one challenge round. \\
    \texttt{giveOpinion} &
    Step 7c &
    Vendor-side answer that updates the binary-search interval. \\
    \texttt{submitCommitment} variants &
    Step 8 &
    Local verification step. Depending on the state, it verifies a generic gate, a left
    boundary condition, or a right/final condition. \\
    \texttt{finalize} or final state transition &
    Step 9 &
    Ends the current dispute execution or applies the protocol consequence determined by the
    sponsor configuration. \\
    \bottomrule
  \end{tabularx}
  \caption{Mapping used to interpret implementation-level gas measurements.}
  \label{tab:methodology-step-mapping}
\end{table}

This mapping also explains why the number of Step 7--8 iterations matters. In normal
sponsored disputes, the protocol may require more than one execution of the challenge and
verification loop because of sponsor rotation. In self-sponsored cases such as \(S_B=B\) and
\(S_V=V\), the same party also plays the relevant sponsor role, so unnecessary repetitions
can be removed. The methodology therefore reports both per-call costs and complete dispute
costs when needed.

\section{Gas accounting scopes}

The most important methodological decision was to define gas scopes before comparing
implementations. The project uses the scopes in Table~\ref{tab:methodology-gas-scopes}.

\begin{table}[htbp]
  \centering
  \begin{tabularx}{\linewidth}{>{\raggedright\arraybackslash}p{0.25\linewidth}>{\raggedright\arraybackslash}X}
    \toprule
    Scope & Meaning \\
    \midrule
    One-time infrastructure &
    Gas paid once to deploy reusable components such as verifier libraries, dispute
    deployers, implementation contracts, or the factory. These costs are not paid again for
    every exchange if the infrastructure is reused. \\
    Marginal optimistic exchange &
    Gas paid for one honest exchange after reusable infrastructure already exists. Depending
    on the variant, this includes clone creation or direct deployment and calls such as
    payment, key release, and completion. \\
    Optimistic path before dispute &
    Gas paid to reach the point where a dispute can be triggered. This includes optimistic
    deployment or creation and the payments or deposits needed before Step 6. \\
    Dispute trigger / initialization &
    Gas paid by \texttt{triggerDispute}. In the measured architecture this may include
    deployment or creation of the dispute contract. \\
    Dispute after trigger &
    Gas paid after the dispute contract exists: challenge rounds, Step 8 verification, and
    finalization. \\
    Single Step 8 call &
    Gas paid by one \texttt{submitCommitment}-type call only. This scope isolates local
    verifier changes, for example the removal of \(\pi_1\) in the hardcoded SHA-256 mode. \\
    End-to-end dispute path &
    Combination of optimistic setup, dispute trigger, post-trigger dispute execution, and
    finalization for a full disputed exchange. \\
    \bottomrule
  \end{tabularx}
  \caption{Gas accounting scopes used in the report.}
  \label{tab:methodology-gas-scopes}
\end{table}

This separation avoids two common mistakes. The first mistake is comparing a one-time
infrastructure deployment to a per-exchange deployment. This would make clone-based
architectures look artificially expensive at setup time or artificially cheap at exchange
time. The second mistake is comparing a single Step 8 optimization with a complete dispute.
For example, removing the generic circuit proof \(\pi_1\) can save gas locally in Step 8,
but the total dispute also includes trigger, challenge rounds, proof verification for other
commitments, gate evaluation, and finalization.

Gas is always taken from transaction receipts, not from informal estimates. In the tests this
is done by waiting for the transaction receipt and reading the \texttt{gasUsed} field. This
matters because storage writes, warm/cold accesses, deployment bytecode, and repeated calls
can make the real gas different from a rough estimate.

\section{Off-chain timing methodology}

Off-chain timings are measured separately from gas. They are used to evaluate the practical
cost of preparing precontracts, computing dispute witnesses, and generating the data needed
for challenge rounds. They do not represent blockchain confirmation time and they are not
converted into gas.

Two execution paths are used:

\begin{itemize}
  \item the WebAssembly path, used by the application and by several tests for small or
  medium inputs;
  \item the native Rust command line path, used for large-file experiments and for the 1 GiB
  hardcoded SHA-256 dispute benchmark.
\end{itemize}

The native path is essential for large files. The browser path is useful for a proof-of-
concept demonstration, but it is not the right benchmark environment for 1 GiB inputs. The
native \path{precontract_cli} measures the cost of reading the input file, encrypting it,
building the circuit, computing commitments, and writing the ciphertext and circuit
artifacts. The native \path{hardcoded_sha256_dispute_cli} measures the off-chain generation
of challenge-round values and the final Step 8 witness data for the hardcoded SHA-256
scenario.

The timing methodology is conservative in interpretation. When a benchmark reports the time
needed to prepare a precontract, this is primarily the vendor-side preparation time in the
current implementation. A fully protocol-level performance table may also account for the
buyer-side verification work, even if the current application does not perform an identical
recomputation before acceptance. Similarly, challenge-response generation is measured for
the party producing the response, while a complete protocol accounting may also include the
counterparty's verification work. The report therefore distinguishes measured implementation
time from conservative protocol-level estimates.

\section{Benchmarking environment}

The on-chain benchmarks are run on a local Hardhat network. The Solidity compiler version is
\texttt{0.8.28}, with optimization enabled, compilation through the intermediate
representation, and optimizer runs set to one. The local network is configured to allow large
contracts, which is useful for testing and for comparing architectures, but it is not a
mainnet guarantee. When mainnet deployability matters, contract bytecode size must be checked
separately.

Most gas tests deploy a mock EntryPoint or local EntryPoint-compatible contract rather than
depending on a live external bundler. This makes the measurements deterministic and keeps the
focus on contract gas rather than on bundler behavior. The interactive application, by
contrast, uses the local Alto bundler and the configured EntryPoint address. Both paths are
useful, but they answer different questions:

\begin{itemize}
  \item Hardhat tests answer: how much gas does this contract path consume?
  \item The full application answers: can the proof of concept be executed through the UI and
  account-abstraction infrastructure?
\end{itemize}

The benchmarks report gas units only. They do not multiply by an ETH price or by a gas price,
because such conversions are time-dependent and would not be stable. They also do not claim
to model L1 congestion, bundler policies, or wallet UX. Those topics are important for
deployment, but the experimental goal of this project is to compare protocol variants under
a controlled local EVM configuration.

\section{Validation strategy}

Validation was performed at several levels.

\paragraph{Functional validation.}
The application was run end-to-end on the local stack. This validated the practical path from
precontract creation to buyer acceptance, contract sponsorship or deployment, buyer payment,
vendor key release, and optimistic completion. This is the path used for demonstrations.

\paragraph{Contract-level validation.}
Hardhat tests exercise the optimistic and dispute contracts directly. They check state
transitions, role restrictions, deployment paths, challenge rounds, and Step 8 submissions.
The performance tests also include assertions on expected values such as the number of
blocks, number of gates, and number of challenge rounds for the generated circuit.

\paragraph{Variant validation.}
Separate tests cover the final specialized architecture, including the normal case,
no \(S\)-deposit, \(S=B\), \(S=V\), self-sponsored dispute execution, and hardcoded
SHA-256. The purpose is not only to obtain gas numbers but also to ensure that the variants
still execute the intended protocol path.

\paragraph{Same-scope comparison.}
Several tests were written specifically to avoid comparing incompatible measurements. For
example, one set of benchmarks measures the current implementation with the same high-level
scope as the previous implementation report, while another set compares monolithic and
specialized variants under the same file size and dispute path. This same-scope discipline
was added because early tables were difficult to interpret when they mixed architectural
versions.

\paragraph{Large-file validation.}
The 1 GiB hardcoded SHA-256 benchmark validates the large-file path with native CLIs. The
test checks that the generated number of blocks and gates matches the expected values,
executes all challenge rounds for the chosen path, submits the final Step 8 witness, and
finalizes the dispute. This validates the feasibility of the large-file pipeline in the
measured scenario.

This validation strategy is practical rather than formal. It increases confidence that the
implementation and measurements correspond to the intended protocol flows, but it does not
replace a full proof of security. In particular, the inherited gate encoding and Merkle
domain-separation questions identified later remain protocol-engineering issues for a future
clean version.
